\chapter{分解理论与多项式图片转写待审核稿}

\begin{dxtips}

本文件为图片与 Markdown 笔记整理出的待审核转写稿，已尽量按 elegantbook 环境归类；正式并入正文前仍建议逐条复核。

\end{dxtips}

\section*{10商域 93 没讲}

\begin{note}[来源上级主题]

第三章 (\%E7\%AC\%AC\%E4\%B8\%89\%E7\%AB\%A0\%201ecd2efa3f4280c7a6b0ea0de416c09f.md)

\end{note}

\begin{note}[10商域 93 没讲]

子级 项目: 定理1没有零因子的交换环R都是一个域Q的子环 (\%E5\%AE\%9A\%E7\%90\%861\%E6\%B2\%A1\%E6\%9C\%89\%E9\%9B\%B6\%E5\%9B\%A0\%E5\%AD\%90\%E7\%9A\%84\%E4\%BA\%A4\%E6\%8D\%A2\%E7\%8E\%AFR\%E9\%83\%BD\%E6\%98\%AF\%E4\%B8\%80\%E4\%B8\%AA\%E5\%9F\%9FQ\%E7\%9A\%84\%E5\%AD\%90\%E7\%8E\%AF\%20202d2efa3f42805194d3ee0dc328bf0c.md), 定理2 Q刚好由所有元组成 (\%E5\%AE\%9A\%E7\%90\%862\%20Q\%E5\%88\%9A\%E5\%A5\%BD\%E7\%94\%B1\%E6\%89\%80\%E6\%9C\%89\%E5\%85\%83\%E7\%BB\%84\%E6\%88\%90\%20202d2efa3f4280869570fa2d2dd1274f.md), 商域的定义 (\%E5\%95\%86\%E5\%9F\%9F\%E7\%9A\%84\%E5\%AE\%9A\%E4\%B9\%89\%20202d2efa3f428055be52faf5725c7e52.md), 定理3 R两个元以上的环  F是包含R的域 F包含R的一个商域 (\%E5\%AE\%9A\%E7\%90\%863\%20R\%E4\%B8\%A4\%E4\%B8\%AA\%E5\%85\%83\%E4\%BB\%A5\%E4\%B8\%8A\%E7\%9A\%84\%E7\%8E\%AF\%20F\%E6\%98\%AF\%E5\%8C\%85\%E5\%90\%ABR\%E7\%9A\%84\%E5\%9F\%9F\%20F\%E5\%8C\%85\%E5\%90\%ABR\%E7\%9A\%84\%E4\%B8\%80\%E4\%B8\%AA\%E5\%95\%86\%E5\%9F\%9F\%20202d2efa3f42802b8611de75dafa7267.md), 定理4同构的环商域也同构 (\%E5\%AE\%9A\%E7\%90\%864\%E5\%90\%8C\%E6\%9E\%84\%E7\%9A\%84\%E7\%8E\%AF\%E5\%95\%86\%E5\%9F\%9F\%E4\%B9\%9F\%E5\%90\%8C\%E6\%9E\%84\%20202d2efa3f42805ea297fb352be96a6e.md)

\end{note}

\begin{note}[来源路径]

近世代数/近世代数/无标题/10商域 93 没讲 1efd2efa3f42806fa86bd0829ae38236.md

\end{note}

\section*{2.6主理想整环与唯一分解整环}

\begin{note}[2.6主理想整环与唯一分解整环]

文本: 59

\end{note}

\begin{note}[来源路径]

近世代数/近世代数/无标题/2 6主理想整环与唯一分解整环 203d2efa3f428041b53bd867ddf75445.md

\end{note}

\section*{2唯一分解环}

\begin{note}[来源上级主题]

第四章整环里面的因子分解 (\%E7\%AC\%AC\%E5\%9B\%9B\%E7\%AB\%A0\%E6\%95\%B4\%E7\%8E\%AF\%E9\%87\%8C\%E9\%9D\%A2\%E7\%9A\%84\%E5\%9B\%A0\%E5\%AD\%90\%E5\%88\%86\%E8\%A7\%A3\%20202d2efa3f42807695c7cb07209ac86f.md)

\end{note}

\begin{note}[2唯一分解环]

待根据图片进一步人工校对。

\end{note}

\begin{note}[来源路径]

近世代数/近世代数/无标题/2唯一分解环 203d2efa3f4280089d71ff9f29c7b591.md

\end{note}

\section*{一元多项式环基础知识}

\begin{note}[来源上级主题]

第四章整环里面的因子分解 (\%E7\%AC\%AC\%E5\%9B\%9B\%E7\%AB\%A0\%E6\%95\%B4\%E7\%8E\%AF\%E9\%87\%8C\%E9\%9D\%A2\%E7\%9A\%84\%E5\%9B\%A0\%E5\%AD\%90\%E5\%88\%86\%E8\%A7\%A3\%20202d2efa3f42807695c7cb07209ac86f.md)

\end{note}

\begin{note}[一元多项式环基础知识]

BATFE ME — MEH I ENTRAR I(r]:

BCR A LF fe BBS

G@) TW itat Ilo JA HS he.

AAT WA iaR Ilo], DR. AW, WR fo) Ile]

BALL

fiem=1 (g@El[z)),

PAHASURRWRKREN, (OM oc) HKRMEFE. RMB f(x),

gael, f(x) BI WH.

f(x) =antayxater+a,x"

\# IlxJW—-t+2UWK. BAF I EME—ARMH, fo) NAM aos ais o>

a, I BA RKAAST.

\end{note}

\begin{note}[图片 OCR 摘录，待人工复核]
以下内容来自源图片识别，便于审核时对照：

\textbf{图片 1}（image.png）
\begin{Verbatim}[fontsize=\small]
BATFE ME — MEH I ENTRAR I(r]:

BCR A LF fe BBS

G@) TW itat Ilo JA HS he.

AAT WA iaR Ilo], DR. AW, WR fo) Ile]
BALL

fiem=1 (g@El[z)),
PAHASURRWRKREN, (OM oc) HKRMEFE. RMB f(x),
gael, f(x) BI WH.
f(x) =antayxater+a,x"

# IlxJW—-t+2UWK. BAF I EME—ARMH, fo) NAM aos ais o>
a, I BA RKAAST.
\end{Verbatim}

\end{note}

\begin{note}[来源路径]

近世代数/近世代数/无标题/一元多项式环基础知识 203d2efa3f428029afb6ff28415761e3.md

\end{note}

\section*{主理想环的定义}

\begin{note}[来源上级主题]

第四章整环里面的因子分解 (\%E7\%AC\%AC\%E5\%9B\%9B\%E7\%AB\%A0\%E6\%95\%B4\%E7\%8E\%AF\%E9\%87\%8C\%E9\%9D\%A2\%E7\%9A\%84\%E5\%9B\%A0\%E5\%AD\%90\%E5\%88\%86\%E8\%A7\%A3\%20202d2efa3f42807695c7cb07209ac86f.md)

\end{note}

\begin{definition}[主理想环的定义]

EM —PEM I SEB, BHT TR Hy

BTU, —-S ERBARIO.

AEB X— > BANG BASS. S| BARS RB.

\end{definition}

\begin{note}[图片 OCR 摘录，待人工复核]
以下内容来自源图片识别，便于审核时对照：

\textbf{图片 1}（image.png）
\begin{Verbatim}[fontsize=\small]
EM —PEM I SEB, BHT TR Hy
BTU, —-S ERBARIO.
AEB X— > BANG BASS. S| BARS RB.
\end{Verbatim}

\end{note}

\begin{note}[来源路径]

近世代数/近世代数/无标题/主理想环的定义 203d2efa3f4280faba6bc969f8b2e4ab.md

\end{note}

\section*{例子}

\begin{note}[来源上级主题]

零因子 (\%E9\%9B\%B6\%E5\%9B\%A0\%E5\%AD\%90\%201ecd2efa3f42800fb79fed095ff2429e.md)

\end{note}

\begin{example}[例子]

fl2 R=(KAthin HMRK). RNS R MeEt—HME

[ajJ+loJ=Lat+o],

FA MIE RMF IXPMERWE MTS. ERMA R ME—THRE.

FET BLE :

La ]Lb]=Lab]                            (2)

Kn HRRKE BRAN SRKA

a=b (modn) 4AM 4n|a—-b it

APR EH. MR

[aJ=[a’], [6J=[0'],

ABZ pa exh

ab—a’b’=a(b—b')+(a—a’)b’

BS Sh iE A

Lab ]=[a’b’].

\end{example}

\begin{note}[图片 OCR 摘录，待人工复核]
以下内容来自源图片识别，便于审核时对照：

\textbf{图片 1}（image.png）
\begin{Verbatim}[fontsize=\small]
fl2 R=(KAthin HMRK). RNS R MeEt—HME
[ajJ+loJ=Lat+o],

FA MIE RMF IXPMERWE MTS. ERMA R ME—THRE.
FET BLE :

La ]Lb]=Lab]                            (2)
Kn HRRKE BRAN SRKA

a=b (modn) 4AM 4n|a—-b it
APR EH. MR
[aJ=[a’], [6J=[0'],
ABZ pa exh
ab—a’b’=a(b—b')+(a—a’)b’

BS Sh iE A

Lab ]=[a’b’].
\end{Verbatim}

\textbf{图片 2}（image 1.png）
\begin{Verbatim}[fontsize=\small]
BRU (2) KET R WAV. LIM MFK WN EMAIL: WKAR E,
FARMS ha. Ale R PER—SH. SAY RE on WRIA.
MR n KERR,
n=ab, nla, n|lb,
HATE R
[aJA[0], (6JAL0], (Ala ]l6]=[ab]=[n]=[0].
AALOERR MZ70, RRR, OAR BAR.
\end{Verbatim}

\end{note}

\begin{note}[来源路径]

近世代数/近世代数/无标题/例子 1ecd2efa3f428040b4fbdcc0b2d26f2a.md

\end{note}

\section*{公因子}

\begin{note}[来源上级主题]

第四章整环里面的因子分解 (\%E7\%AC\%AC\%E5\%9B\%9B\%E7\%AB\%A0\%E6\%95\%B4\%E7\%8E\%AF\%E9\%87\%8C\%E9\%9D\%A2\%E7\%9A\%84\%E5\%9B\%A0\%E5\%AD\%90\%E5\%88\%86\%E8\%A7\%A3\%20202d2efa3f42807695c7cb07209ac86f.md)

\end{note}

\begin{note}[公因子]

EN Foc MMIC a1, a2, +, a, MAAF, Bis c A HE ER a),

Qos "s An.

JE Q15 Az, 5 a, W—-TAAF d Wiilar, az, ws ay HmRABAF,

FEts d BREW Rai, ar, +, a, HBE—-TAAF c BR.

\end{note}

\begin{note}[图片 OCR 摘录，待人工复核]
以下内容来自源图片识别，便于审核时对照：

\textbf{图片 1}（image.png）
\begin{Verbatim}[fontsize=\small]
EN Foc MMIC a1, a2, +, a, MAAF, Bis c A HE ER a),
Qos "s An.

JE Q15 Az, 5 a, W—-TAAF d Wiilar, az, ws ay HmRABAF,
FEts d BREW Rai, ar, +, a, HBE—-TAAF c BR.
\end{Verbatim}

\end{note}

\begin{note}[来源路径]

近世代数/近世代数/无标题/公因子 203d2efa3f4280159f79d1c647f84818.md

\end{note}

\section*{单位 相伴元定义}

\begin{note}[来源上级主题]

第四章整环里面的因子分解 (\%E7\%AC\%AC\%E5\%9B\%9B\%E7\%AB\%A0\%E6\%95\%B4\%E7\%8E\%AF\%E9\%87\%8C\%E9\%9D\%A2\%E7\%9A\%84\%E5\%9B\%A0\%E5\%AD\%90\%E5\%88\%86\%E8\%A7\%A3\%20202d2efa3f42807695c7cb07209ac86f.md)

\end{note}

\begin{definition}[单位 相伴元定义]

单位不是单位元

b是等于a乘一个可逆元 b是a的相伴元

回头算模长的时候别忘了有跟号

单位乘自己逆元是单位元

\end{definition}

\begin{note}[关联图片]
以下图片已纳入本条整理，当前版本优先依据 Markdown 文字整理，图片细节可在审核时对照：

1. image.png\\
2. image 1.png\\
3. image 2.png\\

\end{note}

\begin{note}[来源路径]

近世代数/近世代数/无标题/单位 相伴元定义 202d2efa3f428009a4aec4766781b31c.md

\end{note}

\section*{命题1.2.6 带余除法和阶n的结合}

\begin{note}[来源上级主题]

1.3有限群 (1\%203\%E6\%9C\%89\%E9\%99\%90\%E7\%BE\%A4\%201ead2efa3f4280c39a49d8b0af9aec03.md)

\end{note}

\begin{proposition}[命题1.2.6 带余除法和阶n的结合]

主要要把n拿出来

\end{proposition}

\begin{note}[图片 OCR 摘录，待人工复核]
以下内容来自源图片识别，便于审核时对照：

\textbf{图片 1}（image.png）
\begin{Verbatim}[fontsize=\small]
Aj 1.26

AG = (x) ZAIRE, GE |x| =n, WG = {e,x,x72,--- x7 |}, PPR PA kLAZAAEA

9. AAAS HEY A IRB EY Pe                            “

DETAR AME? WALDL, TEMA PR ORME, ABIL AL O FRSA AT n SRR EMA NTA, FLA nn = |x].
UEW) ATRIA. A, E-SG PRAMAVSRKMAOAMH A MEN BA;s Bo, MOAMN A
ne eA LH

AN RIEHZ-&. ERG PRK x”, Ht me Z. PREPARE, RNA ma=qntr, HH qeZ
O<r<n-1l, MAAAXx" =e, Px axI =x") x =x", MARAT AO ATH An

AM RIEA BOR. RUE, BHRO<m'<m<n-l, fex™ax™, Mx ™ =e, HEL <m-m'<
n-l<n, Pen=|x| LRDHERKK Hx =e, RRSRT TH.

Se LPR, G=f{e,x,x7,--- x "}, SP RB eR EY.
\end{Verbatim}

\end{note}

\begin{note}[来源路径]

近世代数/近世代数/无标题/命题1 2 6 带余除法和阶n的结合 1d2d2efa3f4280f4aba7f41bf31cebda.md

\end{note}

\section*{唯一分解环的定义}

\begin{note}[来源上级主题]

第四章整环里面的因子分解 (\%E7\%AC\%AC\%E5\%9B\%9B\%E7\%AB\%A0\%E6\%95\%B4\%E7\%8E\%AF\%E9\%87\%8C\%E9\%9D\%A2\%E7\%9A\%84\%E5\%9B\%A0\%E5\%AD\%90\%E5\%88\%86\%E8\%A7\%A3\%20202d2efa3f42807695c7cb07209ac86f.md)

\end{note}

\begin{definition}[唯一分解环的定义]

EM —P ER I WY — TSE, HE I WTA S FEA

FE AL ATC RRA ME — A.

EX —-TBRNKB-—B,. —NE—-RM SET AREBEM.

EH 1 — SEAM AA LT ER:

Git) MR—-TAA ATC p HEB BR ab, MBA p EB RHR a Rd.

WEAR Op BERK ad,

\end{definition}

\begin{note}[图片 OCR 摘录，待人工复核]
以下内容来自源图片识别，便于审核时对照：

\textbf{图片 1}（44848cde-04ad-4d01-88d8-c4cd1259715b.png）
\begin{Verbatim}[fontsize=\small]
EM —P ER I WY — TSE, HE I WTA S FEA
FE AL ATC RRA ME — A.

EX —-TBRNKB-—B,. —NE—-RM SET AREBEM.

EH 1 — SEAM AA LT ER:

Git) MR—-TAA ATC p HEB BR ab, MBA p EB RHR a Rd.

WEAR Op BERK ad,
\end{Verbatim}

\end{note}

\begin{note}[来源路径]

近世代数/近世代数/无标题/唯一分解环的定义 203d2efa3f42809a9667cc6b83ca7782.md

\end{note}

\section*{商域的定义}

\begin{note}[来源上级主题]

10商域 93 没讲 (10\%E5\%95\%86\%E5\%9F\%9F\%2093\%20\%E6\%B2\%A1\%E8\%AE\%B2\%201efd2efa3f42806fa86bd0829ae38236.md)

\end{note}

\begin{definition}[商域的定义]

EM —PRQURHRRD—-PRB. BHQBAR, FH QHHHHA

ie-wir

7  (a, bER, 60)

AVE AY.

HERI 2, —SAAAU LHW RASASTHRRASDA-—+t

Fea Sak.

—fHKUL, —PMILAT AERA ERR. RTA

\end{definition}

\begin{note}[图片 OCR 摘录，待人工复核]
以下内容来自源图片识别，便于审核时对照：

\textbf{图片 1}（image.png）
\begin{Verbatim}[fontsize=\small]
EM —PRQURHRRD—-PRB. BHQBAR, FH QHHHHA
ie-wir
7  (a, bER, 60)
AVE AY.
HERI 2, —SAAAU LHW RASASTHRRASDA-—+t
Fea Sak.
—fHKUL, —PMILAT AERA ERR. RTA
\end{Verbatim}

\end{note}

\begin{note}[来源路径]

近世代数/近世代数/无标题/商域的定义 202d2efa3f428055be52faf5725c7e52.md

\end{note}

\section*{因子定义}

\begin{note}[来源上级主题]

第四章整环里面的因子分解 (\%E7\%AC\%AC\%E5\%9B\%9B\%E7\%AB\%A0\%E6\%95\%B4\%E7\%8E\%AF\%E9\%87\%8C\%E9\%9D\%A2\%E7\%9A\%84\%E5\%9B\%A0\%E5\%AD\%90\%E5\%88\%86\%E8\%A7\%A3\%20202d2efa3f42807695c7cb07209ac86f.md)

\end{note}

\begin{definition}[因子定义]

eM Rib. BHI H—-thoce TORI Heb SR, SHE I BRE

Htc, fe

a=be.

fo a AER ER, RRS Ao HAF, FEARS

bla

KHAN. b RB, RNAS

ba

KRM.

FLERBGXPERSMVUE AAR, ICS] AILS RES.

\end{definition}

\begin{note}[图片 OCR 摘录，待人工复核]
以下内容来自源图片识别，便于审核时对照：

\textbf{图片 1}（image.png）
\begin{Verbatim}[fontsize=\small]
eM Rib. BHI H—-thoce TORI Heb SR, SHE I BRE
Htc, fe
a=be.
fo a AER ER, RRS Ao HAF, FEARS
bla
KHAN. b RB, RNAS
ba
KRM.
FLERBGXPERSMVUE AAR, ICS] AILS RES.
\end{Verbatim}

\end{note}

\begin{note}[来源路径]

近世代数/近世代数/无标题/因子定义 202d2efa3f42801594bee964a46cafcd.md

\end{note}

\section*{定义 不可约元}

\begin{note}[来源上级主题]

第四章整环里面的因子分解 (\%E7\%AC\%AC\%E5\%9B\%9B\%E7\%AB\%A0\%E6\%95\%B4\%E7\%8E\%AF\%E9\%87\%8C\%E9\%9D\%A2\%E7\%9A\%84\%E5\%9B\%A0\%E5\%AD\%90\%E5\%88\%86\%E8\%A7\%A3\%20202d2efa3f42807695c7cb07209ac86f.md)

\end{note}

\begin{definition}[定义 不可约元]

单位是有逆元的

证明没有逆元

- 这保证它不是“可逆元素”，否则它的因子分解是平凡的；

- 在 \$\textbackslash{}mathbb\{Z\}\$ 中单位元是 \$\textbackslash{}pm 1\$，在 \$\textbackslash{}mathbb\{Z\}[i]\$ 中是 \$\textbackslash{}pm 1, \textbackslash{}pm i\$ 等。

也就是说：

- 假设存在非单位元素 \$a, b \textbackslash{}in R\$ 使得 p = ab，

- 然后通过反证法或代数结构（如范数、整除性）来推出矛盾，说明不可能；

- 所以只能有一个因子是单位元，即 p 是不可约的。

\end{definition}

\begin{note}[关联图片]
以下图片已纳入本条整理，当前版本优先依据 Markdown 文字整理，图片细节可在审核时对照：

1. image.png\\

\end{note}

\begin{note}[来源路径]

近世代数/近世代数/无标题/定义 不可约元 202d2efa3f428089be02ea14d2704267.md

\end{note}

\section*{定义2.18整环没有零因子是交换环}

\begin{note}[来源上级主题]

环 (\%E7\%8E\%AF\%201d9d2efa3f4280f38e94de87f834c519.md)

\end{note}

\begin{definition}[定义2.18整环没有零因子是交换环]

特别是类加群的时候，凡事不和阶互诉的元都是零因子

\end{definition}

\begin{note}[关联图片]
以下图片已纳入本条整理，当前版本优先依据 Markdown 文字整理，图片细节可在审核时对照：

1. image.png\\

\end{note}

\begin{note}[来源路径]

近世代数/近世代数/无标题/定义2 18整环没有零因子是交换环 1d8d2efa3f428099bbfed4ab53ac0f69.md

\end{note}

\section*{定义2.25主理想整环}

\begin{definition}[定义2.25主理想整环]

每一个理想都是主理想

\end{definition}

\begin{note}[图片 OCR 摘录，待人工复核]
以下内容来自源图片识别，便于审核时对照：

\textbf{图片 1}（image.png）
\begin{Verbatim}[fontsize=\small]
A PEAR aS ee ERE EA, BT DA Re cE a FT PE
EFUES

x SL 2.25
A(R, +,°) RP RAH, WAT R RDS EPR, SRAE-PBEH, MAA—-PBMIIR PRE
F270, PPT UAE M

I=(a)=Ra                                                      (O22)

MIX.                                                                                                                               .

ESAS B HY EITAR GS, (EPA AI. Bi SL EE EB (Zz, +, -) 0

Pia, Fee UEHY (Z, +, -) xe-S EB.
\end{Verbatim}

\end{note}

\begin{note}[来源路径]

近世代数/近世代数/无标题/定义2 25主理想整环 203d2efa3f4280079f4dc6da3c39ca7c.md

\end{note}

\section*{定义唯一分解}

\begin{note}[来源上级主题]

第四章整环里面的因子分解 (\%E7\%AC\%AC\%E5\%9B\%9B\%E7\%AB\%A0\%E6\%95\%B4\%E7\%8E\%AF\%E9\%87\%8C\%E9\%9D\%A2\%E7\%9A\%84\%E5\%9B\%A0\%E5\%AD\%90\%E5\%88\%86\%E8\%A7\%A3\%20202d2efa3f42807695c7cb07209ac86f.md)

\end{note}

\begin{definition}[定义唯一分解]

eM Rib, —TH\#HI H—-ts0e £1 BAM—DR, HUTA

BE Bk is E :

G)         a=pipop, (6; HI WADA);

Gi) An5R AT

a=nggs (qi INABA),

BBA

r=s.

FFARTA UE a: WK R— FP, et

gi=eip; (eo: ZI WE.

\end{definition}

\begin{note}[图片 OCR 摘录，待人工复核]
以下内容来自源图片识别，便于审核时对照：

\textbf{图片 1}（image.png）
\begin{Verbatim}[fontsize=\small]
eM Rib, —TH#HI H—-ts0e £1 BAM—DR, HUTA
BE Bk is E :
G)         a=pipop, (6; HI WADA);
Gi) An5R AT
a=nggs (qi INABA),
BBA
r=s.
FFARTA UE a: WK R— FP, et
gi=eip; (eo: ZI WE.
\end{Verbatim}

\end{note}

\begin{note}[来源路径]

近世代数/近世代数/无标题/定义唯一分解 202d2efa3f42804fbb82c9b5883516db.md

\end{note}

\section*{定义平凡因子与真因子}

\begin{note}[来源上级主题]

第四章整环里面的因子分解 (\%E7\%AC\%AC\%E5\%9B\%9B\%E7\%AB\%A0\%E6\%95\%B4\%E7\%8E\%AF\%E9\%87\%8C\%E9\%9D\%A2\%E7\%9A\%84\%E5\%9B\%A0\%E5\%AD\%90\%E5\%88\%86\%E8\%A7\%A3\%20202d2efa3f42807695c7cb07209ac86f.md)

\end{note}

\begin{definition}[定义平凡因子与真因子]

相伴元环单位是评分因子

逆元是乘法的

\end{definition}

\begin{note}[图片 OCR 摘录，待人工复核]
以下内容来自源图片识别，便于审核时对照：

\textbf{图片 1}（image.png）
\begin{Verbatim}[fontsize=\small]
eM BUR HACK HERAT. HAD HAT. Ra
RAW, Wie HAF.

REILRNB-B, —PHBRM p AETAEM. —TbRM p HARE
of Xt 7S BE BE ER, AA 1 Mop RYU pv. ART ROSA
SBR p RAAF. KR ELAM RE, +1 MEERA, p=lp,
—p=(—-Dp Ep MHA. FURNTWK., Re pH-TPERE, CF
AAENLAT. RM p RAA—-PMER, BE PAO R11 KRRAH ee
PE it BNA
\end{Verbatim}

\end{note}

\begin{note}[来源路径]

近世代数/近世代数/无标题/定义平凡因子与真因子 202d2efa3f42803fa3b6ff5c998de7b7.md

\end{note}

\section*{定理 主理想环I是一个唯一分解环}

\begin{note}[来源上级主题]

第四章整环里面的因子分解 (\%E7\%AC\%AC\%E5\%9B\%9B\%E7\%AB\%A0\%E6\%95\%B4\%E7\%8E\%AF\%E9\%87\%8C\%E9\%9D\%A2\%E7\%9A\%84\%E5\%9B\%A0\%E5\%AD\%90\%E5\%88\%86\%E8\%A7\%A3\%20202d2efa3f42807695c7cb07209ac86f.md)

\end{note}

\begin{theorem}[定理 主理想环I是一个唯一分解环]

待根据图片进一步人工校对。

\end{theorem}

\begin{note}[来源路径]

近世代数/近世代数/无标题/定理 主理想环I是一个唯一分解环 203d2efa3f4280f1bf1ac75fef055fb8.md

\end{note}

\section*{定理1 I是唯一分解环I[x]也是}

\begin{note}[来源上级主题]

第四章整环里面的因子分解 (\%E7\%AC\%AC\%E5\%9B\%9B\%E7\%AB\%A0\%E6\%95\%B4\%E7\%8E\%AF\%E9\%87\%8C\%E9\%9D\%A2\%E7\%9A\%84\%E5\%9B\%A0\%E5\%AD\%90\%E5\%88\%86\%E8\%A7\%A3\%20202d2efa3f42807695c7cb07209ac86f.md)

\end{note}

\begin{theorem}[定理1 I是唯一分解环I[x]也是]

待根据图片进一步人工校对。

\end{theorem}

\begin{note}[来源路径]

近世代数/近世代数/无标题/定理1 I是唯一分解环I[x]也是 203d2efa3f42800fa68fd8427dfa8780.md

\end{note}

\section*{定理1 单位乘单位还是单位 单位的逆还是单位}

\begin{note}[来源上级主题]

第四章整环里面的因子分解 (\%E7\%AC\%AC\%E5\%9B\%9B\%E7\%AB\%A0\%E6\%95\%B4\%E7\%8E\%AF\%E9\%87\%8C\%E9\%9D\%A2\%E7\%9A\%84\%E5\%9B\%A0\%E5\%AD\%90\%E5\%88\%86\%E8\%A7\%A3\%20202d2efa3f42807695c7cb07209ac86f.md)

\end{note}

\begin{theorem}[定理1 单位乘单位还是单位 单位的逆还是单位]

1] AM ite Me MRR Het. Bite Mitte th

A fi.

\end{theorem}

\begin{note}[图片 OCR 摘录，待人工复核]
以下内容来自源图片识别，便于审核时对照：

\textbf{图片 1}（image.png）
\begin{Verbatim}[fontsize=\small]
1] AM ite Me MRR Het. Bite Mitte th
A fi.
\end{Verbatim}

\end{note}

\begin{note}[来源路径]

近世代数/近世代数/无标题/定理1 单位乘单位还是单位 单位的逆还是单位 202d2efa3f4280028d54e2faca1b169d.md

\end{note}

\section*{定理1唯一分解环性质}

\begin{note}[来源上级主题]

第四章整环里面的因子分解 (\%E7\%AC\%AC\%E5\%9B\%9B\%E7\%AB\%A0\%E6\%95\%B4\%E7\%8E\%AF\%E9\%87\%8C\%E9\%9D\%A2\%E7\%9A\%84\%E5\%9B\%A0\%E5\%AD\%90\%E5\%88\%86\%E8\%A7\%A3\%20202d2efa3f42807695c7cb07209ac86f.md)

\end{note}

\begin{theorem}[定理1唯一分解环性质]

EEL —S EA RIN LP HE

Git) MR—-TAA AT p HEBER ab, BA p BRR a Md.

\end{theorem}

\begin{note}[图片 OCR 摘录，待人工复核]
以下内容来自源图片识别，便于审核时对照：

\textbf{图片 1}（image.png）
\begin{Verbatim}[fontsize=\small]
EEL —S EA RIN LP HE
Git) MR—-TAA AT p HEBER ab, BA p BRR a Md.
\end{Verbatim}

\end{note}

\begin{note}[来源路径]

近世代数/近世代数/无标题/定理1唯一分解环性质 203d2efa3f4280a88b26de5685fbb8de.md

\end{note}

\section*{定理1无零因子环所有不为0的元对加法阶都是一样的}

\begin{note}[来源上级主题]

第三章 (\%E7\%AC\%AC\%E4\%B8\%89\%E7\%AB\%A0\%201ecd2efa3f4280c7a6b0ea0de416c09f.md)

\end{note}

\begin{theorem}[定理1无零因子环所有不为0的元对加法阶都是一样的]

子级 项目: 证明 (\%E8\%AF\%81\%E6\%98\%8E\%201edd2efa3f42808b9ca2d57ff1844f03.md)

\end{theorem}

\begin{note}[关联图片]
以下图片已纳入本条整理，当前版本优先依据 Markdown 文字整理，图片细节可在审核时对照：

1. image.png\\

\end{note}

\begin{note}[来源路径]

近世代数/近世代数/无标题/定理1无零因子环所有不为0的元对加法阶都是一样的 1edd2efa3f428098a9d1d9282007041a.md

\end{note}

\section*{定理1欧式环是一个主理想环，是一个唯一分解环}

\begin{note}[来源上级主题]

第四章整环里面的因子分解 (\%E7\%AC\%AC\%E5\%9B\%9B\%E7\%AB\%A0\%E6\%95\%B4\%E7\%8E\%AF\%E9\%87\%8C\%E9\%9D\%A2\%E7\%9A\%84\%E5\%9B\%A0\%E5\%AD\%90\%E5\%88\%86\%E8\%A7\%A3\%20202d2efa3f42807695c7cb07209ac86f.md)

\end{note}

\begin{theorem}[定理1欧式环是一个主理想环，是一个唯一分解环]

EH 1 (EMR KRH I —ER-TFEBRH, Ail—ER— SEAM.

\end{theorem}

\begin{note}[图片 OCR 摘录，待人工复核]
以下内容来自源图片识别，便于审核时对照：

\textbf{图片 1}（image.png）
\begin{Verbatim}[fontsize=\small]
EH 1 (EMR KRH I —ER-TFEBRH, Ail—ER— SEAM.
\end{Verbatim}

\end{note}

\begin{note}[来源路径]

近世代数/近世代数/无标题/定理1欧式环是一个主理想环，是一个唯一分解环 203d2efa3f42809a9537c3fbfb56d4bf.md

\end{note}

\section*{定理1没有零因子的交换环R都是一个域Q的子环}

\begin{note}[来源上级主题]

10商域 93 没讲 (10\%E5\%95\%86\%E5\%9F\%9F\%2093\%20\%E6\%B2\%A1\%E8\%AE\%B2\%201efd2efa3f42806fa86bd0829ae38236.md)

\end{note}

\begin{theorem}[定理1没有零因子的交换环R都是一个域Q的子环]

ER 1 B—TRASAF WKH R ME-TMQWFEH.

\end{theorem}

\begin{note}[图片 OCR 摘录，待人工复核]
以下内容来自源图片识别，便于审核时对照：

\textbf{图片 1}（image.png）
\begin{Verbatim}[fontsize=\small]
ER 1 B—TRASAF WKH R ME-TMQWFEH.
\end{Verbatim}

\end{note}

\begin{note}[来源路径]

近世代数/近世代数/无标题/定理1没有零因子的交换环R都是一个域Q的子环 202d2efa3f42805194d3ee0dc328bf0c.md

\end{note}

\section*{定理2 I是唯一分解环还有什么是}

\begin{note}[来源上级主题]

第四章整环里面的因子分解 (\%E7\%AC\%AC\%E5\%9B\%9B\%E7\%AB\%A0\%E6\%95\%B4\%E7\%8E\%AF\%E9\%87\%8C\%E9\%9D\%A2\%E7\%9A\%84\%E5\%9B\%A0\%E5\%AD\%90\%E5\%88\%86\%E8\%A7\%A3\%20202d2efa3f42807695c7cb07209ac86f.md)

\end{note}

\begin{theorem}[定理2 I是唯一分解环还有什么是]

EH2 WRI BMH, BAT Lx, r2., °°, c, JHE, KB x,

Koy, Ly ET EW DRA ET.

HEH 1, ST NT, Ilo SER. (RAT AG,

BURAK A ETERS BHR S7 3). KH, RNA T—THE-aH

PAE ERRAW AF.

\end{theorem}

\begin{note}[图片 OCR 摘录，待人工复核]
以下内容来自源图片识别，便于审核时对照：

\textbf{图片 1}（image.png）
\begin{Verbatim}[fontsize=\small]
EH2 WRI BMH, BAT Lx, r2., °°, c, JHE, KB x,
Koy, Ly ET EW DRA ET.

HEH 1, ST NT, Ilo SER. (RAT AG,
BURAK A ETERS BHR S7 3). KH, RNA T—THE-aH
PAE ERRAW AF.
\end{Verbatim}

\end{note}

\begin{note}[来源路径]

近世代数/近世代数/无标题/定理2 I是唯一分解环还有什么是 203d2efa3f42804f8335ecd40e8ebe42.md

\end{note}

\section*{定理2 Q刚好由所有元组成}

\begin{note}[来源上级主题]

10商域 93 没讲 (10\%E5\%95\%86\%E5\%9F\%9F\%2093\%20\%E6\%B2\%A1\%E8\%AE\%B2\%201efd2efa3f42806fa86bd0829ae38236.md)

\end{note}

\begin{theorem}[定理2 Q刚好由所有元组成]

EH2 QUA AMAT

5 (a, BER, 640)

PRE BOA, ix

Fm ab ba,

\end{theorem}

\begin{note}[图片 OCR 摘录，待人工复核]
以下内容来自源图片识别，便于审核时对照：

\textbf{图片 1}（image.png）
\begin{Verbatim}[fontsize=\small]
EH2 QUA AMAT
5 (a, BER, 640)
PRE BOA, ix
Fm ab ba,
\end{Verbatim}

\end{note}

\begin{note}[来源路径]

近世代数/近世代数/无标题/定理2 Q刚好由所有元组成 202d2efa3f4280869570fa2d2dd1274f.md

\end{note}

\section*{定理2 整环需要什么变成唯一分解环}

\begin{note}[来源上级主题]

第四章整环里面的因子分解 (\%E7\%AC\%AC\%E5\%9B\%9B\%E7\%AB\%A0\%E6\%95\%B4\%E7\%8E\%AF\%E9\%87\%8C\%E9\%9D\%A2\%E7\%9A\%84\%E5\%9B\%A0\%E5\%AD\%90\%E5\%88\%86\%E8\%A7\%A3\%20202d2efa3f42807695c7cb07209ac86f.md)

\end{note}

\begin{theorem}[定理2 整环需要什么变成唯一分解环]

不可约元都是素元

不是零元和单位的元都有一个分解

\end{theorem}

\begin{note}[图片 OCR 摘录，待人工复核]
以下内容来自源图片识别，便于审核时对照：

\textbf{图片 1}（image.png）
\begin{Verbatim}[fontsize=\small]
EH 2 BE—-TMEMR I AU PER:
@ 1 BB—-TRRBE BREAN TE a BATS
a=Ppipop, (pi BIWAW AI);
Git) I MAA ATC ERIC
BERS I ERE — SE — 5h EH.
\end{Verbatim}

\end{note}

\begin{note}[来源路径]

近世代数/近世代数/无标题/定理2 整环需要什么变成唯一分解环 203d2efa3f42804598d3f020bf976baf.md

\end{note}

\section*{定理2单位乘不可约元还是不可约元}

\begin{note}[来源上级主题]

第四章整环里面的因子分解 (\%E7\%AC\%AC\%E5\%9B\%9B\%E7\%AB\%A0\%E6\%95\%B4\%E7\%8E\%AF\%E9\%87\%8C\%E9\%9D\%A2\%E7\%9A\%84\%E5\%9B\%A0\%E5\%AD\%90\%E5\%88\%86\%E8\%A7\%A3\%20202d2efa3f42807695c7cb07209ac86f.md)

\end{note}

\begin{theorem}[定理2单位乘不可约元还是不可约元]

EH2 Bie WADA SHR ep He—-PAM AT.

HERA eH Fe AO, PAO, MBMRASAT, MU epA0. ep HASH

BAL, ASR

1=e’(ep)=(e’e) p.

pH, SREARG.

DLE b ep WAT, FH DAB. BA

ep=bc, p=ble'c), b| p.

\{A p BARA AT, b RRB, Alt b ER p WAP:

b=e"p=(e'e (ep) (eM).

MRE, b ep WHIT, ANE 1, ce EMM. KR ep RAFN

AF. ue SE.

\end{theorem}

\begin{note}[图片 OCR 摘录，待人工复核]
以下内容来自源图片识别，便于审核时对照：

\textbf{图片 1}（image.png）
\begin{Verbatim}[fontsize=\small]
EH2 Bie WADA SHR ep He—-PAM AT.

HERA eH Fe AO, PAO, MBMRASAT, MU epA0. ep HASH
BAL, ASR

1=e’(ep)=(e’e) p.

pH, SREARG.

DLE b ep WAT, FH DAB. BA

ep=bc, p=ble'c), b| p.
{A p BARA AT, b RRB, Alt b ER p WAP:
b=e"p=(e'e (ep) (eM).

MRE, b ep WHIT, ANE 1, ce EMM. KR ep RAFN
AF. ue SE.
\end{Verbatim}

\end{note}

\begin{note}[来源路径]

近世代数/近世代数/无标题/定理2单位乘不可约元还是不可约元 202d2efa3f4280c8893ee029f5f5168d.md

\end{note}

\section*{定理2整数环是一个主理想环}

\begin{note}[来源上级主题]

第四章整环里面的因子分解 (\%E7\%AC\%AC\%E5\%9B\%9B\%E7\%AB\%A0\%E6\%95\%B4\%E7\%8E\%AF\%E9\%87\%8C\%E9\%9D\%A2\%E7\%9A\%84\%E5\%9B\%A0\%E5\%AD\%90\%E5\%88\%86\%E8\%A7\%A3\%20202d2efa3f42807695c7cb07209ac86f.md)

\end{note}

\begin{theorem}[定理2整数环是一个主理想环]

EB2 BRHE—-SEBRH, Al e—ThHE—AMBH.

Fa — FS SLA RR RPM EHS. BFE

\end{theorem}

\begin{note}[图片 OCR 摘录，待人工复核]
以下内容来自源图片识别，便于审核时对照：

\textbf{图片 1}（image.png）
\begin{Verbatim}[fontsize=\small]
EB2 BRHE—-SEBRH, Al e—ThHE—AMBH.
Fa — FS SLA RR RPM EHS. BFE
\end{Verbatim}

\end{note}

\begin{note}[来源路径]

近世代数/近世代数/无标题/定理2整数环是一个主理想环 203d2efa3f4280cba00fc015db3e2793.md

\end{note}

\section*{定理3 R两个元以上的环 F是包含R的域 F包含R的一个商域}

\begin{note}[来源上级主题]

10商域 93 没讲 (10\%E5\%95\%86\%E5\%9F\%9F\%2093\%20\%E6\%B2\%A1\%E8\%AE\%B2\%201efd2efa3f42806fa86bd0829ae38236.md)

\end{note}

\begin{theorem}[定理3 R两个元以上的环 F是包含R的域 F包含R的一个商域]

3 RER BE-+AATUENTNS, FE—-TF+MSR WR. HA

F @GR W—S Fm.

WEAR EF

ab =b a=  (a, bER, b¥0)

ABM. EF MF

O= \{sree\}  (a, BER, b¥0).

Q BRER H—-TAM. IES.

\{LR OR—PBRMBAAHTASMMD., MiHBMMDCeRRF R KH

MEMFVE; RRB. RK HPRNWECAERRAFR He. MURNA

\end{theorem}

\begin{note}[图片 OCR 摘录，待人工复核]
以下内容来自源图片识别，便于审核时对照：

\textbf{图片 1}（image.png）
\begin{Verbatim}[fontsize=\small]
3 RER BE-+AATUENTNS, FE—-TF+MSR WR. HA
F @GR W—S Fm.
WEAR EF
ab =b a=  (a, bER, b¥0)
ABM. EF MF
O= {sree}  (a, BER, b¥0).
Q BRER H—-TAM. IES.
{LR OR—PBRMBAAHTASMMD., MiHBMMDCeRRF R KH
MEMFVE; RRB. RK HPRNWECAERRAFR He. MURNA
\end{Verbatim}

\end{note}

\begin{note}[来源路径]

近世代数/近世代数/无标题/定理3 R两个元以上的环 F是包含R的域 F包含R的一个商域 202d2efa3f42802b8611de75dafa7267.md

\end{note}

\section*{定理3 整环不等元零的元有真因子}

\begin{note}[来源上级主题]

第四章整环里面的因子分解 (\%E7\%AC\%AC\%E5\%9B\%9B\%E7\%AB\%A0\%E6\%95\%B4\%E7\%8E\%AF\%E9\%87\%8C\%E9\%9D\%A2\%E7\%9A\%84\%E5\%9B\%A0\%E5\%AD\%90\%E5\%88\%86\%E8\%A7\%A3\%20202d2efa3f42807695c7cb07209ac86f.md)

\end{note}

\begin{theorem}[定理3 整环不等元零的元有真因子]

EH3 BHP—-TASTSH a ARRAS HRERER

a=be,

b Alc MAGMA.

\end{theorem}

\begin{note}[图片 OCR 摘录，待人工复核]
以下内容来自源图片识别，便于审核时对照：

\textbf{图片 1}（image.png）
\begin{Verbatim}[fontsize=\small]
EH3 BHP—-TASTSH a ARRAS HRERER
a=be,
b Alc MAGMA.
\end{Verbatim}

\end{note}

\begin{note}[来源路径]

近世代数/近世代数/无标题/定理3 整环不等元零的元有真因子 202d2efa3f4280529816da6994b32844.md

\end{note}

\section*{定理3 最大公因子与单位因子}

\begin{note}[来源上级主题]

第四章整环里面的因子分解 (\%E7\%AC\%AC\%E5\%9B\%9B\%E7\%AB\%A0\%E6\%95\%B4\%E7\%8E\%AF\%E9\%87\%8C\%E9\%9D\%A2\%E7\%9A\%84\%E5\%9B\%A0\%E5\%AD\%90\%E5\%88\%86\%E8\%A7\%A3\%20202d2efa3f42807695c7cb07209ac86f.md)

\end{note}

\begin{theorem}[定理3 最大公因子与单位因子]

3 —SE—ARH I Wh ca Mb EI B-CARKAAF. a

Alb HIT RAKAAT d Ald’ RiH-T+PMAF:

d’=ed (ce \#M).

\end{theorem}

\begin{note}[图片 OCR 摘录，待人工复核]
以下内容来自源图片识别，便于审核时对照：

\textbf{图片 1}（image.png）
\begin{Verbatim}[fontsize=\small]
3 —SE—ARH I Wh ca Mb EI B-CARKAAF. a
Alb HIT RAKAAT d Ald’ RiH-T+PMAF:
d’=ed (ce #M).
\end{Verbatim}

\end{note}

\begin{note}[来源路径]

近世代数/近世代数/无标题/定理3 最大公因子与单位因子 203d2efa3f428022b1aefef6d6ccbfc7.md

\end{note}

\section*{定理3一元多项式环都是欧式环}

\begin{note}[来源上级主题]

第四章整环里面的因子分解 (\%E7\%AC\%AC\%E5\%9B\%9B\%E7\%AB\%A0\%E6\%95\%B4\%E7\%8E\%AF\%E9\%87\%8C\%E9\%9D\%A2\%E7\%9A\%84\%E5\%9B\%A0\%E5\%AD\%90\%E5\%88\%86\%E8\%A7\%A3\%20202d2efa3f42807695c7cb07209ac86f.md)

\end{note}

\begin{theorem}[定理3一元多项式环都是欧式环]

EMS —PSRF LMW—TEMRH Flr ]E—TMRH.

\end{theorem}

\begin{note}[图片 OCR 摘录，待人工复核]
以下内容来自源图片识别，便于审核时对照：

\textbf{图片 1}（image.png）
\begin{Verbatim}[fontsize=\small]
EMS —PSRF LMW—TEMRH Flr ]E—TMRH.
\end{Verbatim}

\end{note}

\begin{note}[来源路径]

近世代数/近世代数/无标题/定理3一元多项式环都是欧式环 203d2efa3f428054b7ece00218586ac4.md

\end{note}

\section*{定理3有限群的阶都是父群因子的阶}

\begin{note}[来源上级主题]

子群的陪集 (\%E5\%AD\%90\%E7\%BE\%A4\%E7\%9A\%84\%E9\%99\%AA\%E9\%9B\%86\%201c3d2efa3f42806a85b0f0113028b207.md)

\end{note}

\begin{theorem}[定理3有限群的阶都是父群因子的阶]

N=nj.                        bE 5G.

EH 3 —SARHG HE—titc HRI n HERG HH.

WEARS saa AE MP OT. Le. on BRC AM. GER.

\end{theorem}

\begin{note}[图片 OCR 摘录，待人工复核]
以下内容来自源图片识别，便于审核时对照：

\textbf{图片 1}（image.png）
\begin{Verbatim}[fontsize=\small]
N=nj.                        bE 5G.
EH 3 —SARHG HE—titc HRI n HERG HH.
WEARS saa AE MP OT. Le. on BRC AM. GER.
\end{Verbatim}

\end{note}

\begin{note}[来源路径]

近世代数/近世代数/无标题/定理3有限群的阶都是父群因子的阶 1c3d2efa3f4280c2992cc6ef4f5024c8.md

\end{note}

\section*{定理4同构的环商域也同构}

\begin{note}[来源上级主题]

10商域 93 没讲 (10\%E5\%95\%86\%E5\%9F\%9F\%2093\%20\%E6\%B2\%A1\%E8\%AE\%B2\%201efd2efa3f42806fa86bd0829ae38236.md)

\end{note}

\begin{theorem}[定理4同构的环商域也同构]

一个环最多只有一个商域

\end{theorem}

\begin{note}[图片 OCR 摘录，待人工复核]
以下内容来自源图片识别，便于审核时对照：

\textbf{图片 1}（image.png）
\begin{Verbatim}[fontsize=\small]
EH4 AMHHAH HMA. RHMAMRKA, —TbHRSAA-T
96 e
\end{Verbatim}

\end{note}

\begin{note}[来源路径]

近世代数/近世代数/无标题/定理4同构的环商域也同构 202d2efa3f42805ea297fb352be96a6e.md

\end{note}

\section*{左零因子}

\begin{note}[来源上级主题]

零因子 (\%E9\%9B\%B6\%E5\%9B\%A0\%E5\%AD\%90\%201ecd2efa3f42800fb79fed095ff2429e.md)

\end{note}

\begin{note}[左零因子]

交换环左零因子就是右，非交换环不一定

\end{note}

\begin{note}[图片 OCR 摘录，待人工复核]
以下内容来自源图片识别，便于审核时对照：

\textbf{图片 1}（image.png）
\begin{Verbatim}[fontsize=\small]
eM WRE-THE
a%0, bA0, (KH ab=0,
KTR, a ZXTHH-TESATF, 6#-TASAFT. A. AFATFH
MESA F.
\end{Verbatim}

\end{note}

\begin{note}[来源路径]

近世代数/近世代数/无标题/左零因子 1ecd2efa3f428041bb3bcc3b3ca95daa.md

\end{note}

\section*{引理1 有限序列判定}

\begin{note}[来源上级主题]

第四章整环里面的因子分解 (\%E7\%AC\%AC\%E5\%9B\%9B\%E7\%AB\%A0\%E6\%95\%B4\%E7\%8E\%AF\%E9\%87\%8C\%E9\%9D\%A2\%E7\%9A\%84\%E5\%9B\%A0\%E5\%AD\%90\%E5\%88\%86\%E8\%A7\%A3\%20202d2efa3f42807695c7cb07209ac86f.md)

\end{note}

\begin{lemma}[引理1 有限序列判定]

S132 1 (eI BP ERM. WERE

Qi» Aes a3, 7° (a;ET)

BEATA HTS, WARNER —EE—TH ESI.

\end{lemma}

\begin{note}[图片 OCR 摘录，待人工复核]
以下内容来自源图片识别，便于审核时对照：

\textbf{图片 1}（image.png）
\begin{Verbatim}[fontsize=\small]
S132 1 (eI BP ERM. WERE
Qi» Aes a3, 7° (a;ET)
BEATA HTS, WARNER —EE—TH ESI.
\end{Verbatim}

\end{note}

\begin{note}[来源路径]

近世代数/近世代数/无标题/引理1 有限序列判定 203d2efa3f4280bea27cd23fb527b89e.md

\end{note}

\section*{引理1 本原多项式只能游本源多项式相乘得到}

\begin{note}[来源上级主题]

第四章整环里面的因子分解 (\%E7\%AC\%AC\%E5\%9B\%9B\%E7\%AB\%A0\%E6\%95\%B4\%E7\%8E\%AF\%E9\%87\%8C\%E9\%9D\%A2\%E7\%9A\%84\%E5\%9B\%A0\%E5\%AD\%90\%E5\%88\%86\%E8\%A7\%A3\%20202d2efa3f42807695c7cb07209ac86f.md)

\end{note}

\begin{lemma}[引理1 本原多项式只能游本源多项式相乘得到]

BIH 1 fe f(2=—gladh(c). RA f(oeARSHRN, SAMY

g(x) Al A (x) MEA RS KAT.

\end{lemma}

\begin{note}[图片 OCR 摘录，待人工复核]
以下内容来自源图片识别，便于审核时对照：

\textbf{图片 1}（image.png）
\begin{Verbatim}[fontsize=\small]
BIH 1 fe f(2=—gladh(c). RA f(oeARSHRN, SAMY
g(x) Al A (x) MEA RS KAT.
\end{Verbatim}

\end{note}

\begin{note}[来源路径]

近世代数/近世代数/无标题/引理1 本原多项式只能游本源多项式相乘得到 203d2efa3f42804a99cef3d1a3008763.md

\end{note}

\section*{引理2 主理想环的一个不可约元生成一个极大理想}

\begin{note}[来源上级主题]

第四章整环里面的因子分解 (\%E7\%AC\%AC\%E5\%9B\%9B\%E7\%AB\%A0\%E6\%95\%B4\%E7\%8E\%AF\%E9\%87\%8C\%E9\%9D\%A2\%E7\%9A\%84\%E5\%9B\%A0\%E5\%AD\%90\%E5\%88\%86\%E8\%A7\%A3\%20202d2efa3f42807695c7cb07209ac86f.md)

\end{note}

\begin{lemma}[引理2 主理想环的一个不可约元生成一个极大理想]

待根据图片进一步人工校对。

\end{lemma}

\begin{note}[来源路径]

近世代数/近世代数/无标题/引理2 主理想环的一个不可约元生成一个极大理想 203d2efa3f42806c99eac511dd8df111.md

\end{note}

\section*{引理2 给不等于0的多项式改形式}

\begin{note}[来源上级主题]

第四章整环里面的因子分解 (\%E7\%AC\%AC\%E5\%9B\%9B\%E7\%AB\%A0\%E6\%95\%B4\%E7\%8E\%AF\%E9\%87\%8C\%E9\%9D\%A2\%E7\%9A\%84\%E5\%9B\%A0\%E5\%AD\%90\%E5\%88\%86\%E8\%A7\%A3\%20202d2efa3f42807695c7cb07209ac86f.md)

\end{note}

\begin{lemma}[引理2 给不等于0的多项式改形式]

5122 QlreIWMB—-*+ASFEHZHARK f(OMAUS KH

fla) =2 fo)

ABs, xX Ha, FEI, f(y I[c | WRRBSMWR. MR go (ewF fox)

WE, BA

Sol(xad=efo(x) Ce HI WFD.

\end{lemma}

\begin{note}[图片 OCR 摘录，待人工复核]
以下内容来自源图片识别，便于审核时对照：

\textbf{图片 1}（image.png）
\begin{Verbatim}[fontsize=\small]
5122 QlreIWMB—-*+ASFEHZHARK f(OMAUS KH
fla) =2 fo)
ABs, xX Ha, FEI, f(y I[c | WRRBSMWR. MR go (ewF fox)
WE, BA
Sol(xad=efo(x) Ce HI WFD.
\end{Verbatim}

\end{note}

\begin{note}[来源路径]

近世代数/近世代数/无标题/引理2 给不等于0的多项式改形式 203d2efa3f42801981f3fb6c052e7210.md

\end{note}

\section*{引理3 I[x]可约充要条件就是Q[x]可约}

\begin{note}[来源上级主题]

第四章整环里面的因子分解 (\%E7\%AC\%AC\%E5\%9B\%9B\%E7\%AB\%A0\%E6\%95\%B4\%E7\%8E\%AF\%E9\%87\%8C\%E9\%9D\%A2\%E7\%9A\%84\%E5\%9B\%A0\%E5\%AD\%90\%E5\%88\%86\%E8\%A7\%A3\%20202d2efa3f42807695c7cb07209ac86f.md)

\end{note}

\begin{lemma}[引理3 I[x]可约充要条件就是Q[x]可约]

B13 Iz ]M—*+ARS MR f. (oc) Ie BONER

fox) \# Qle BAA.

\end{lemma}

\begin{note}[图片 OCR 摘录，待人工复核]
以下内容来自源图片识别，便于审核时对照：

\textbf{图片 1}（image.png）
\begin{Verbatim}[fontsize=\small]
B13 Iz ]M—*+ARS MR f. (oc) Ie BONER
fox) # Qle BAA.
\end{Verbatim}

\end{note}

\begin{note}[来源路径]

近世代数/近世代数/无标题/引理3 I[x]可约充要条件就是Q[x]可约 203d2efa3f428041b03adc66ef2f8383.md

\end{note}

\section*{引理4次数大于0的本原多项式有唯一分解}

\begin{note}[来源上级主题]

第四章整环里面的因子分解 (\%E7\%AC\%AC\%E5\%9B\%9B\%E7\%AB\%A0\%E6\%95\%B4\%E7\%8E\%AF\%E9\%87\%8C\%E9\%9D\%A2\%E7\%9A\%84\%E5\%9B\%A0\%E5\%AD\%90\%E5\%88\%86\%E8\%A7\%A3\%20202d2efa3f42807695c7cb07209ac86f.md)

\end{note}

\begin{lemma}[引理4次数大于0的本原多项式有唯一分解]

on 4 Ix |WH—-*+RAKFEWARS AK fo (x) Ix] BAM

\end{lemma}

\begin{note}[图片 OCR 摘录，待人工复核]
以下内容来自源图片识别，便于审核时对照：

\textbf{图片 1}（image.png）
\begin{Verbatim}[fontsize=\small]
on 4 Ix |WH—-*+RAKFEWARS AK fo (x) Ix] BAM
\end{Verbatim}

\end{note}

\begin{note}[来源路径]

近世代数/近世代数/无标题/引理4次数大于0的本原多项式有唯一分解 203d2efa3f42803984b2de025e50e146.md

\end{note}

\section*{引理多项式环都是主理想环}

\begin{note}[来源上级主题]

第四章整环里面的因子分解 (\%E7\%AC\%AC\%E5\%9B\%9B\%E7\%AB\%A0\%E6\%95\%B4\%E7\%8E\%AF\%E9\%87\%8C\%E9\%9D\%A2\%E7\%9A\%84\%E5\%9B\%A0\%E5\%AD\%90\%E5\%88\%86\%E8\%A7\%A3\%20202d2efa3f42807695c7cb07209ac86f.md)

\end{note}

\begin{lemma}[引理多项式环都是主理想环]

512 Be Iz ]ABH I ELW—AeHRH, Ie] 7G

g(x) =a,zr"t+a,-12" ' ++

Wm KRARM oc, ZI N-TbB. BA Ice IPERSMHAK f(x) MAW

Bm

f(x=qlarglxa)tr(xe) (q(x), r(x) E€lILex])

WIR, Br (2) =0 Rr (ce) WRBDF 2 (oc) WU nn.

\end{lemma}

\begin{note}[图片 OCR 摘录，待人工复核]
以下内容来自源图片识别，便于审核时对照：

\textbf{图片 1}（image.png）
\begin{Verbatim}[fontsize=\small]
512 Be Iz ]ABH I ELW—AeHRH, Ie] 7G
g(x) =a,zr"t+a,-12" ' ++
Wm KRARM oc, ZI N-TbB. BA Ice IPERSMHAK f(x) MAW
Bm
f(x=qlarglxa)tr(xe) (q(x), r(x) E€lILex])
WIR, Br (2) =0 Rr (ce) WRBDF 2 (oc) WU nn.
\end{Verbatim}

\end{note}

\begin{note}[来源路径]

近世代数/近世代数/无标题/引理多项式环都是主理想环 203d2efa3f4280e387c9e6c5e5952fd6.md

\end{note}

\section*{推论}

\begin{note}[来源上级主题]

零因子与消去律互相证明 (\%E9\%9B\%B6\%E5\%9B\%A0\%E5\%AD\%90\%E4\%B8\%8E\%E6\%B6\%88\%E5\%8E\%BB\%E5\%BE\%8B\%E4\%BA\%92\%E7\%9B\%B8\%E8\%AF\%81\%E6\%98\%8E\%201ecd2efa3f42809fa11dc863d390a502.md)

\end{note}

\begin{note}[推论]

Hit A-THEWMRA-THEERM, MAA-TRER ERY.

LAER TUVGR T-TREE AAS = FRIAR Es SSE RA I ACH

ff. BOT EPMUNAE, BETSRSATHAEE. —TbH AMAIA

WARY EMMA, WRAL E= PIMA eal S.

\end{note}

\begin{note}[图片 OCR 摘录，待人工复核]
以下内容来自源图片识别，便于审核时对照：

\textbf{图片 1}（image.png）
\begin{Verbatim}[fontsize=\small]
Hit A-THEWMRA-THEERM, MAA-TRER ERY.

LAER TUVGR T-TREE AAS = FRIAR Es SSE RA I ACH
ff. BOT EPMUNAE, BETSRSATHAEE. —TbH AMAIA
WARY EMMA, WRAL E= PIMA eal S.
\end{Verbatim}

\textbf{图片 2}（image 1.png）
\begin{Verbatim}[fontsize=\small]
=, SSZMMBBXR (BAM)

eA                      AEB ME B

AUT = HKE EAE

WAR = RASAF EAE

AUT = RASAF EAE (WASP)
RASAF = WA WIAA IZ, HH
\end{Verbatim}

\end{note}

\begin{note}[来源路径]

近世代数/近世代数/无标题/推论 1ecd2efa3f4280a6bd0ac80610cbea45.md

\end{note}

\section*{推论 互素的定义}

\begin{note}[来源上级主题]

第四章整环里面的因子分解 (\%E7\%AC\%AC\%E5\%9B\%9B\%E7\%AB\%A0\%E6\%95\%B4\%E7\%8E\%AF\%E9\%87\%8C\%E9\%9D\%A2\%E7\%9A\%84\%E5\%9B\%A0\%E5\%AD\%90\%E5\%88\%86\%E8\%A7\%A3\%20202d2efa3f42807695c7cb07209ac86f.md)

\end{note}

\begin{definition}[推论 互素的定义]

ie — TEAR IT Wn SoGa1> ae,» a, EI B-BARKE

AF. a1, a2, > a, WPATRACZA FT RETR WA T.

, ARILPTMR—-TRALAF ETH, BULTIC EA —*

RAKLAF HET EL. FAK, RMT SER BL

EDR — TBS.

EM RNG, —TS EAR A WIC a1, a2, , a, BR, RHE

BKAAT 2B Mi.

\end{definition}

\begin{note}[图片 OCR 摘录，待人工复核]
以下内容来自源图片识别，便于审核时对照：

\textbf{图片 1}（image.png）
\begin{Verbatim}[fontsize=\small]
ie — TEAR IT Wn SoGa1> ae,» a, EI B-BARKE
AF. a1, a2, > a, WPATRACZA FT RETR WA T.

, ARILPTMR—-TRALAF ETH, BULTIC EA —*
RAKLAF HET EL. FAK, RMT SER BL
EDR — TBS.

EM RNG, —TS EAR A WIC a1, a2, , a, BR, RHE
BKAAT 2B Mi.
\end{Verbatim}

\end{note}

\begin{note}[来源路径]

近世代数/近世代数/无标题/推论 互素的定义 203d2efa3f4280c8a6bbeab8b7b50fd4.md

\end{note}

\section*{无标题}

\begin{note}[来源上级主题]

第四章整环里面的因子分解 (\%E7\%AC\%AC\%E5\%9B\%9B\%E7\%AB\%A0\%E6\%95\%B4\%E7\%8E\%AF\%E9\%87\%8C\%E9\%9D\%A2\%E7\%9A\%84\%E5\%9B\%A0\%E5\%AD\%90\%E5\%88\%86\%E8\%A7\%A3\%20202d2efa3f42807695c7cb07209ac86f.md)

\end{note}

\begin{note}[无标题]

待根据图片进一步人工校对。

\end{note}

\begin{note}[来源路径]

近世代数/近世代数/无标题/无标题 204d2efa3f4280a9a1a2fad9e464a923.md

\end{note}

\section*{无零因子作用}

\begin{note}[来源上级主题]

3.3 (3\%203\%20209d2efa3f428020b15ef2f886bfb11e.md)

\end{note}

\begin{note}[无零因子作用]

待根据图片进一步人工校对。

\end{note}

\begin{note}[来源路径]

近世代数/近世代数/无标题/无零因子作用 20ad2efa3f428015bfb2db4f02684706.md

\end{note}

\section*{本原多项式定义}

\begin{note}[来源上级主题]

第四章整环里面的因子分解 (\%E7\%AC\%AC\%E5\%9B\%9B\%E7\%AB\%A0\%E6\%95\%B4\%E7\%8E\%AF\%E9\%87\%8C\%E9\%9D\%A2\%E7\%9A\%84\%E5\%9B\%A0\%E5\%AD\%90\%E5\%88\%86\%E8\%A7\%A3\%20202d2efa3f42807695c7cb07209ac86f.md)

\end{note}

\begin{definition}[本原多项式定义]

HEM Ile l]W—+h50 fo) wm—-TARRSHAK, BH fo HARN

KABT ii.

RIXPEM, BR

Gi) —TARASWAARSE FS.

Gi) MRARBSMK f(xo)yA, BA

f(o=glayh(x),

MB g(a) A oNKRAKTES, AMBDT fo WRK

\end{definition}

\begin{note}[图片 OCR 摘录，待人工复核]
以下内容来自源图片识别，便于审核时对照：

\textbf{图片 1}（image.png）
\begin{Verbatim}[fontsize=\small]
HEM Ile l]W—+h50 fo) wm—-TARRSHAK, BH fo HARN
KABT ii.
RIXPEM, BR
Gi) —TARASWAARSE FS.
Gi) MRARBSMK f(xo)yA, BA
f(o=glayh(x),
MB g(a) A oNKRAKTES, AMBDT fo WRK
\end{Verbatim}

\end{note}

\begin{note}[来源路径]

近世代数/近世代数/无标题/本原多项式定义 203d2efa3f4280808b6be7296615d881.md

\end{note}

\section*{欧式环定义}

\begin{note}[来源上级主题]

第四章整环里面的因子分解 (\%E7\%AC\%AC\%E5\%9B\%9B\%E7\%AB\%A0\%E6\%95\%B4\%E7\%8E\%AF\%E9\%87\%8C\%E9\%9D\%A2\%E7\%9A\%84\%E5\%9B\%A0\%E5\%AD\%90\%E5\%88\%86\%E8\%A7\%A3\%20202d2efa3f42807695c7cb07209ac86f.md)

\end{note}

\begin{definition}[欧式环定义]

欧式环

\end{definition}

\begin{note}[图片 OCR 摘录，待人工复核]
以下内容来自源图片识别，便于审核时对照：

\textbf{图片 1}（image.png）
\begin{Verbatim}[fontsize=\small]
EM HPS RE I MY — PRR, 2298
Ll. A—TM I WIEST ERR A BSE ARR @ FEE
Il. ET I W—-*+RSFEH7Ca, I HER SITUS
b=qatr (q, r€l)
Wes, x Br=0 x g(r)<g(a).
\end{Verbatim}

\end{note}

\begin{note}[来源路径]

近世代数/近世代数/无标题/欧式环定义 203d2efa3f4280ea8dd9fc99cd6ba87e.md

\end{note}

\section*{环域主理想之间的关系}

\begin{note}[来源上级主题]

第四章整环里面的因子分解 (\%E7\%AC\%AC\%E5\%9B\%9B\%E7\%AB\%A0\%E6\%95\%B4\%E7\%8E\%AF\%E9\%87\%8C\%E9\%9D\%A2\%E7\%9A\%84\%E5\%9B\%A0\%E5\%AD\%90\%E5\%88\%86\%E8\%A7\%A3\%20202d2efa3f42807695c7cb07209ac86f.md)

\end{note}

\begin{note}[环域主理想之间的关系]

panko fen                                BOR

an

+f \{0\} \#324 > WKIBAA ]

\end{note}

\begin{note}[图片 OCR 摘录，待人工复核]
以下内容来自源图片识别，便于审核时对照：

\textbf{图片 1}（image.png）
\begin{Verbatim}[fontsize=\small]
panko fen                                BOR
an
+f {0} #324 > WKIBAA ]
\end{Verbatim}

\end{note}

\begin{note}[来源路径]

近世代数/近世代数/无标题/环域主理想之间的关系 204d2efa3f4280a9a86fd593f2874f38.md

\end{note}

\section*{第四章整环里面的因子分解}

\begin{note}[第四章整环里面的因子分解]

子级 项目: 因子定义 (\%E5\%9B\%A0\%E5\%AD\%90\%E5\%AE\%9A\%E4\%B9\%89\%20202d2efa3f42801594bee964a46cafcd.md), 环域主理想之间的关系 (\%E7\%8E\%AF\%E5\%9F\%9F\%E4\%B8\%BB\%E7\%90\%86\%E6\%83\%B3\%E4\%B9\%8B\%E9\%97\%B4\%E7\%9A\%84\%E5\%85\%B3\%E7\%B3\%BB\%20204d2efa3f4280a9a86fd593f2874f38.md), 单位 相伴元定义 (\%E5\%8D\%95\%E4\%BD\%8D\%20\%E7\%9B\%B8\%E4\%BC\%B4\%E5\%85\%83\%E5\%AE\%9A\%E4\%B9\%89\%20202d2efa3f428009a4aec4766781b31c.md), 定理1 单位乘单位还是单位 单位的逆还是单位 (\%E5\%AE\%9A\%E7\%90\%861\%20\%E5\%8D\%95\%E4\%BD\%8D\%E4\%B9\%98\%E5\%8D\%95\%E4\%BD\%8D\%E8\%BF\%98\%E6\%98\%AF\%E5\%8D\%95\%E4\%BD\%8D\%20\%E5\%8D\%95\%E4\%BD\%8D\%E7\%9A\%84\%E9\%80\%86\%E8\%BF\%98\%E6\%98\%AF\%E5\%8D\%95\%E4\%BD\%8D\%20202d2efa3f4280028d54e2faca1b169d.md), 定义平凡因子与真因子 (\%E5\%AE\%9A\%E4\%B9\%89\%E5\%B9\%B3\%E5\%87\%A1\%E5\%9B\%A0\%E5\%AD\%90\%E4\%B8\%8E\%E7\%9C\%9F\%E5\%9B\%A0\%E5\%AD\%90\%20202d2efa3f42803fa3b6ff5c998de7b7.md), 定义 不可约元 (\%E5\%AE\%9A\%E4\%B9\%89\%20\%E4\%B8\%8D\%E5\%8F\%AF\%E7\%BA\%A6\%E5\%85\%83\%20202d2efa3f428089be02ea14d2704267.md), 定理2单位乘不可约元还是不可约元 (\%E5\%AE\%9A\%E7\%90\%862\%E5\%8D\%95\%E4\%BD\%8D\%E4\%B9\%98\%E4\%B8\%8D\%E5\%8F\%AF\%E7\%BA\%A6\%E5\%85\%83\%E8\%BF\%98\%E6\%98\%AF\%E4\%B8\%8D\%E5\%8F\%AF\%E7\%BA\%A6\%E5\%85\%83\%20202d2efa3f4280c8893ee029f5f5168d.md), 定理3 整环不等元零的元有真因子 (\%E5\%AE\%9A\%E7\%90\%863\%20\%E6\%95\%B4\%E7\%8E\%AF\%E4\%B8\%8D\%E7\%AD\%89\%E5\%85\%83\%E9\%9B\%B6\%E7\%9A\%84\%E5\%85\%83\%E6\%9C\%89\%E7\%9C\%9F\%E5\%9B\%A0\%E5\%AD\%90\%20202d2efa3f4280529816da6994b32844.md), 定义唯一分解 (\%E5\%AE\%9A\%E4\%B9\%89\%E5\%94\%AF\%E4\%B8\%80\%E5\%88\%86\%E8\%A7\%A3\%20202d2efa3f42804fbb82c9b5883516db.md), 2唯一分解环 (2\%E5\%94\%AF\%E4\%B8\%80\%E5\%88\%86\%E8\%A7\%A3\%E7\%8E\%AF\%20203d2efa3f4280089d71ff9f29c7b591.md), 唯一分解环的定义 (\%E5\%94\%AF\%E4\%B8\%80\%E5\%88\%86\%E8\%A7\%A3\%E7\%8E\%AF\%E7\%9A\%84\%E5\%AE\%9A\%E4\%B9\%89\%20203d2efa3f42809a9667cc6b83ca7782.md), 定理1唯一分解环性质 (\%E5\%AE\%9A\%E7\%90\%861\%E5\%94\%AF\%E4\%B8\%80\%E5\%88\%86\%E8\%A7\%A3\%E7\%8E\%AF\%E6\%80\%A7\%E8\%B4\%A8\%20203d2efa3f4280a88b26de5685fbb8de.md), 定理2 整环需要什么变成唯一分解环 (\%E5\%AE\%9A\%E7\%90\%862\%20\%E6\%95\%B4\%E7\%8E\%AF\%E9\%9C\%80\%E8\%A6\%81\%E4\%BB\%80\%E4\%B9\%88\%E5\%8F\%98\%E6\%88\%90\%E5\%94\%AF\%E4\%B8\%80\%E5\%88\%86\%E8\%A7\%A3\%E7\%8E\%AF\%20203d2efa3f42804598d3f020bf976baf.md), 公因子 (\%E5\%85\%AC\%E5\%9B\%A0\%E5\%AD\%90\%20203d2efa3f4280159f79d1c647f84818.md), 定理3 最大公因子与单位因子 (\%E5\%AE\%9A\%E7\%90\%863\%20\%E6\%9C\%80\%E5\%A4\%A7\%E5\%85\%AC\%E5\%9B\%A0\%E5\%AD\%90\%E4\%B8\%8E\%E5\%8D\%95\%E4\%BD\%8D\%E5\%9B\%A0\%E5\%AD\%90\%20203d2efa3f428022b1aefef6d6ccbfc7.md), 推论 互素的定义 (\%E6\%8E\%A8\%E8\%AE\%BA\%20\%E4\%BA\%92\%E7\%B4\%A0\%E7\%9A\%84\%E5\%AE\%9A\%E4\%B9\%89\%20203d2efa3f4280c8a6bbeab8b7b50fd4.md), 主理想环的定义 (\%E4\%B8\%BB\%E7\%90\%86\%E6\%83\%B3\%E7\%8E\%AF\%E7\%9A\%84\%E5\%AE\%9A\%E4\%B9\%89\%20203d2efa3f4280faba6bc969f8b2e4ab.md), 引理1 有限序列判定 (\%E5\%BC\%95\%E7\%90\%861\%20\%E6\%9C\%89\%E9\%99\%90\%E5\%BA\%8F\%E5\%88\%97\%E5\%88\%A4\%E5\%AE\%9A\%20203d2efa3f4280bea27cd23fb527b89e.md), 引理2 主理想环的一个不可约元生成一个极大理想 (\%E5\%BC\%95\%E7\%90\%862\%20\%E4\%B8\%BB\%E7\%90\%86\%E6\%83\%B3\%E7\%8E\%AF\%E7\%9A\%84\%E4\%B8\%80\%E4\%B8\%AA\%E4\%B8\%8D\%E5\%8F\%AF\%E7\%BA\%A6\%E5\%85\%83\%E7\%94\%9F\%E6\%88\%90\%E4\%B8\%80\%E4\%B8\%AA\%E6\%9E\%81\%E5\%A4\%A7\%E7\%90\%86\%E6\%83\%B3\%20203d2efa3f42806c99eac511dd8df111.md), 定理 主理想环I是一个唯一分解环 (\%E5\%AE\%9A\%E7\%90\%86\%20\%E4\%B8\%BB\%E7\%90\%86\%E6\%83\%B3\%E7\%8E\%AFI\%E6\%98\%AF\%E4\%B8\%80\%E4\%B8\%AA\%E5\%94\%AF\%E4\%B8\%80\%E5\%88\%86\%E8\%A7\%A3\%E7\%8E\%AF\%20203d2efa3f4280f1bf1ac75fef055fb8.md), 欧式环定义 (\%E6\%AC\%A7\%E5\%BC\%8F\%E7\%8E\%AF\%E5\%AE\%9A\%E4\%B9\%89\%20203d2efa3f4280ea8dd9fc99cd6ba87e.md), 定理1欧式环是一个主理想环，是一个唯一分解环 (\%E5\%AE\%9A\%E7\%90\%861\%E6\%AC\%A7\%E5\%BC\%8F\%E7\%8E\%AF\%E6\%98\%AF\%E4\%B8\%80\%E4\%B8\%AA\%E4\%B8\%BB\%E7\%90\%86\%E6\%83\%B3\%E7\%8E\%AF\%EF\%BC\%8C\%E6\%98\%AF\%E4\%B8\%80\%E4\%B8\%AA\%E5\%94\%AF\%E4\%B8\%80\%E5\%88\%86\%E8\%A7\%A3\%E7\%8E\%AF\%20203d2efa3f42809a9537c3fbfb56d4bf.md), 定理2整数环是一个主理想环 (\%E5\%AE\%9A\%E7\%90\%862\%E6\%95\%B4\%E6\%95\%B0\%E7\%8E\%AF\%E6\%98\%AF\%E4\%B8\%80\%E4\%B8\%AA\%E4\%B8\%BB\%E7\%90\%86\%E6\%83\%B3\%E7\%8E\%AF\%20203d2efa3f4280cba00fc015db3e2793.md), 引理多项式环都是主理想环 (\%E5\%BC\%95\%E7\%90\%86\%E5\%A4\%9A\%E9\%A1\%B9\%E5\%BC\%8F\%E7\%8E\%AF\%E9\%83\%BD\%E6\%98\%AF\%E4\%B8\%BB\%E7\%90\%86\%E6\%83\%B3\%E7\%8E\%AF\%20203d2efa3f4280e387c9e6c5e5952fd6.md), 定理3一元多项式环都是欧式环 (\%E5\%AE\%9A\%E7\%90\%863\%E4\%B8\%80\%E5\%85\%83\%E5\%A4\%9A\%E9\%A1\%B9\%E5\%BC\%8F\%E7\%8E\%AF\%E9\%83\%BD\%E6\%98\%AF\%E6\%AC\%A7\%E5\%BC\%8F\%E7\%8E\%AF\%20203d2efa3f428054b7ece00218586ac4.md), 一元多项式环基础知识 (\%E4\%B8\%80\%E5\%85\%83\%E5\%A4\%9A\%E9\%A1\%B9\%E5\%BC\%8F\%E7\%8E\%AF\%E5\%9F\%BA\%E7\%A1\%80\%E7\%9F\%A5\%E8\%AF\%86\%20203d2efa3f428029afb6ff28415761e3.md), 本原多项式定义 (\%E6\%9C\%AC\%E5\%8E\%9F\%E5\%A4\%9A\%E9\%A1\%B9\%E5\%BC\%8F\%E5\%AE\%9A\%E4\%B9\%89\%20203d2efa3f4280808b6be7296615d881.md), 引理1 本原多项式只能游本源多项式相乘得到 (\%E5\%BC\%95\%E7\%90\%861\%20\%E6\%9C\%AC\%E5\%8E\%9F\%E5\%A4\%9A\%E9\%A1\%B9\%E5\%BC\%8F\%E5\%8F\%AA\%E8\%83\%BD\%E6\%B8\%B8\%E6\%9C\%AC\%E6\%BA\%90\%E5\%A4\%9A\%E9\%A1\%B9\%E5\%BC\%8F\%E7\%9B\%B8\%E4\%B9\%98\%E5\%BE\%97\%E5\%88\%B0\%20203d2efa3f42804a99cef3d1a3008763.md), 引理2 给不等于0的多项式改形式 (\%E5\%BC\%95\%E7\%90\%862\%20\%E7\%BB\%99\%E4\%B8\%8D\%E7\%AD\%89\%E4\%BA\%8E0\%E7\%9A\%84\%E5\%A4\%9A\%E9\%A1\%B9\%E5\%BC\%8F\%E6\%94\%B9\%E5\%BD\%A2\%E5\%BC\%8F\%20203d2efa3f42801981f3fb6c052e7210.md), 引理3 I[x]可约充要条件就是Q[x]可约 (\%E5\%BC\%95\%E7\%90\%863\%20I\%5Bx\%5D\%E5\%8F\%AF\%E7\%BA\%A6\%E5\%85\%85\%E8\%A6\%81\%E6\%9D\%A1\%E4\%BB\%B6\%E5\%B0\%B1\%E6\%98\%AFQ\%5Bx\%5D\%E5\%8F\%AF\%E7\%BA\%A6\%20203d2efa3f428041b03adc66ef2f8383.md), 引理4次数大于0的本原多项式有唯一分解 (\%E5\%BC\%95\%E7\%90\%864\%E6\%AC\%A1\%E6\%95\%B0\%E5\%A4\%A7\%E4\%BA\%8E0\%E7\%9A\%84\%E6\%9C\%AC\%E5\%8E\%9F\%E5\%A4\%9A\%E9\%A1\%B9\%E5\%BC\%8F\%E6\%9C\%89\%E5\%94\%AF\%E4\%B8\%80\%E5\%88\%86\%E8\%A7\%A3\%20203d2efa3f42803984b2de025e50e146.md), 定理1 I是唯一分解环I[x]也是 (\%E5\%AE\%9A\%E7\%90\%861\%20I\%E6\%98\%AF\%E5\%94\%AF\%E4\%B8\%80\%E5\%88\%86\%E8\%A7\%A3\%E7\%8E\%AFI\%5Bx\%5D\%E4\%B9\%9F\%E6\%98\%AF\%20203d2efa3f42800fa68fd8427dfa8780.md), 定理2 I是唯一分解环还有什么是 (\%E5\%AE\%9A\%E7\%90\%862\%20I\%E6\%98\%AF\%E5\%94\%AF\%E4\%B8\%80\%E5\%88\%86\%E8\%A7\%A3\%E7\%8E\%AF\%E8\%BF\%98\%E6\%9C\%89\%E4\%BB\%80\%E4\%B9\%88\%E6\%98\%AF\%20203d2efa3f42804f8335ecd40e8ebe42.md), 无标题 (\%E6\%97\%A0\%E6\%A0\%87\%E9\%A2\%98\%20204d2efa3f4280a9a1a2fad9e464a923.md)

\end{note}

\begin{note}[来源路径]

近世代数/近世代数/无标题/第四章整环里面的因子分解 202d2efa3f42807695c7cb07209ac86f.md

\end{note}

\section*{证明}

\begin{note}[来源上级主题]

定理1无零因子环所有不为0的元对加法阶都是一样的 (\%E5\%AE\%9A\%E7\%90\%861\%E6\%97\%A0\%E9\%9B\%B6\%E5\%9B\%A0\%E5\%AD\%90\%E7\%8E\%AF\%E6\%89\%80\%E6\%9C\%89\%E4\%B8\%8D\%E4\%B8\%BA0\%E7\%9A\%84\%E5\%85\%83\%E5\%AF\%B9\%E5\%8A\%A0\%E6\%B3\%95\%E9\%98\%B6\%E9\%83\%BD\%E6\%98\%AF\%E4\%B8\%80\%E6\%A0\%B7\%E7\%9A\%84\%201edd2efa3f428098a9d1d9282007041a.md)

\end{note}

\begin{proof}

这里针对的都是加法因为没有零因子只有n个a相加是0和n个b相加是零两种可能性

\end{proof}

\begin{note}[关联图片]
以下图片已纳入本条整理，当前版本优先依据 Markdown 文字整理，图片细节可在审核时对照：

1. image.png\\

\end{note}

\begin{note}[来源路径]

近世代数/近世代数/无标题/证明 1edd2efa3f42808b9ca2d57ff1844f03.md

\end{note}

\section*{零因子}

\begin{note}[来源上级主题]

第三章 (\%E7\%AC\%AC\%E4\%B8\%89\%E7\%AB\%A0\%201ecd2efa3f4280c7a6b0ea0de416c09f.md)

\end{note}

\begin{note}[零因子]

子级 项目: 例子 (\%E4\%BE\%8B\%E5\%AD\%90\%201ecd2efa3f428040b4fbdcc0b2d26f2a.md), 左零因子 (\%E5\%B7\%A6\%E9\%9B\%B6\%E5\%9B\%A0\%E5\%AD\%90\%201ecd2efa3f428041bb3bcc3b3ca95daa.md)

\end{note}

\begin{note}[关联图片]
以下图片已纳入本条整理，当前版本优先依据 Markdown 文字整理，图片细节可在审核时对照：

1. image.png\\

\end{note}

\begin{note}[来源路径]

近世代数/近世代数/无标题/零因子 1ecd2efa3f42800fb79fed095ff2429e.md

\end{note}

\section*{零因子与消去律互相证明}

\begin{note}[来源上级主题]

第三章 (\%E7\%AC\%AC\%E4\%B8\%89\%E7\%AB\%A0\%201ecd2efa3f4280c7a6b0ea0de416c09f.md)

\end{note}

\begin{proof}

子级 项目: 推论 (\%E6\%8E\%A8\%E8\%AE\%BA\%201ecd2efa3f4280a6bd0ac80610cbea45.md)

\end{proof}

\begin{note}[来源路径]

近世代数/近世代数/无标题/零因子与消去律互相证明 1ecd2efa3f42809fa11dc863d390a502.md

\end{note}

\section*{零因子作用}

\begin{note}[零因子作用]

SA THERES ION INA BBE LA.

WEAR: @RICKR, BAHIA].

\#7 | a |= m,ma = 0, Vb ¢ 0,(ma)b = a(mb) = 0

=> mb = 0 => |b| < |a|

fa] HY 484 [B| = [a| =>] B [= mn.

EX: PASAT IIE CATAL

CRNA VLE TAR AK SME

\end{note}

\begin{note}[图片 OCR 摘录，待人工复核]
以下内容来自源图片识别，便于审核时对照：

\textbf{图片 1}（image.png）
\begin{Verbatim}[fontsize=\small]
SA THERES ION INA BBE LA.
WEAR: @RICKR, BAHIA].
#7 | a |= m,ma = 0, Vb ¢ 0,(ma)b = a(mb) = 0
=> mb = 0 => |b| < |a|
fa] HY 484 [B| = [a| =>] B [= mn.
EX: PASAT IIE CATAL
CRNA VLE TAR AK SME
\end{Verbatim}

\end{note}

\begin{note}[来源路径]

近世代数/近世代数/无标题/零因子作用 20ed2efa3f428023a13dc8880e5e0a69.md

\end{note}
